\section{Related Work}
\label{sec:related}

Catastrophic forgetting (CF), first characterized by~\cite{mccloskey1989catastrophic}, remains a
fundamental obstacle to continual learning in neural networks: sequential gradient updates cause a
model to overwrite previously acquired knowledge when it adapts to a new task distribution.
With the rise of large language models (LLMs), CF has taken on renewed urgency; practitioners
increasingly need to fine-tune pre-trained models on successive domain- or task-specific corpora
without retaining all prior training data~\cite{luo2023empirical,li2024revisiting}.
We organize the existing literature into four clusters.

% Required packages (add to main.tex preamble if missing):
%   \usepackage{booktabs}

\subsection{Regularization-Based Methods}

Regularization approaches penalize changes to model parameters that are deemed important for prior
tasks.
Elastic Weight Consolidation (EWC)~\cite{kirkpatrick2017overcoming} anchors each parameter to its
post-training value with a quadratic penalty weighted by the diagonal of the empirical Fisher
information matrix, providing a theoretically grounded stability--plasticity trade-off.
While EWC is widely adopted as a baseline, its memory cost grows linearly with the number of tasks
when the full Fisher diagonal must be stored, and the penalty strength $\lambda$ must be tuned
per task sequence.
Hierarchical regularization~\cite{song2025hierarchical} extends the EWC idea to transformer layers
by assigning task-specific importance scores at different network depths, improving plasticity while
preserving stability in cross-domain continual instruction tuning.
Sharpness-aware minimization (SAM) has also been explored as an implicit regularizer that encourages
convergence to flat minima, which generalize better under distribution shift, though its compute
overhead is substantial.
These methods share a common limitation: mispecification of the regularization strength leads to
either underfitting (high stability, low plasticity) or catastrophic forgetting (low stability,
high plasticity).

\subsection{Rehearsal-Based Methods}

Rehearsal approaches maintain a small episodic memory of prior-task examples and interleave them
with current-task training.
Experience Replay (ER) is the simplest instantiation~\cite{luo2023empirical}: a fixed-size buffer
is populated uniformly from past data, and buffer samples are concatenated with each mini-batch.
Structured Selective Replay (SSR)~\cite{huang2024ssr} improves upon na\"{i}ve ER by selecting which
examples to retain based on gradient-space diversity, substantially reducing the buffer size required
to achieve a given retention level.
Mixed training~\cite{reynolds2025mixed} relaxes the strict sequential assumption by jointly sampling
from the current task and all prior-task corpora at each step, effectively treating continual
learning as multi-task learning with delayed data availability.
Rehearsal methods are appealing because they impose no architectural constraints, but retaining prior
data raises privacy and storage concerns in deployment contexts subject to data-use regulations.

\subsection{Architectural and PEFT-Based Methods}

Parameter-efficient fine-tuning (PEFT) methods reduce the number of trainable parameters, which
naturally limits interference between tasks.
LoRA~\cite{hu2021lora} constrains weight updates to a low-rank factorization
$W + \Delta W = W + AB$, updating only $A \in \mathbb{R}^{d \times r}$ and
$B \in \mathbb{R}^{r \times k}$ while freezing $W$.
Orthogonal LoRA (O-LoRA)~\cite{wang2023olora} strengthens this by imposing an orthogonality
constraint between the LoRA subspace learned for the current task and those of prior tasks,
directly minimizing interference in parameter space without storing any training examples.
CURLoRA~\cite{fawi2024curlora} applies CUR matrix decomposition to initialize adapter weights from
prior tasks, leveraging structured low-rank approximations to reduce forgetting during sequential
fine-tuning.
N-LoRA~\cite{yang2024nlora} incorporates a null-space projection mechanism that updates parameters
only in directions orthogonal to the gradient subspace of prior tasks, enforcing forward-transfer
compatibility.
GainLoRA~\cite{liang2025gainlora} introduces a gain-scaling reparameterization that modulates the
effective rank of each adapter layer dynamically based on task-specific curvature estimates.
SECURA~\cite{zhang2025secura} proposes a security-aware rehearsal framework that augments LoRA
adapters with curriculum-guided memory consolidation to address both forgetting and adversarial
robustness.
A comprehensive comparative study~\cite{haque2025comparative} surveys many of these adapter-based
continual learning strategies across diverse benchmarks, providing a useful reference for
method selection in practice.

\subsection{Compute-Efficient and Small-Model Studies}

Despite the abundance of continual learning methods, most empirical studies are conducted at scale
(7B+ parameters) with ample GPU resources, often training on high-memory A100 or H100 hardware.
The interaction between quantization, low-rank adaptation, and forgetting at the 1--4B parameter
scale remains poorly characterized.
QLoRA~\cite{dettmers2023qlora} demonstrated that 4-bit NF4 quantization with LoRA adapters trained
in \texttt{bfloat16} matches full-precision fine-tuning quality on \emph{static} tasks, but its
behavior under sequential task shifts has not been systematically evaluated.
Post-training continual adaptation~\cite{kumar2025posttraining} examines lightweight fine-tuning
strategies for already-instruction-tuned models but focuses on single-task adaptation rather than
multi-task sequences.
A broader survey~\cite{nayak2025sculpting} explores sculpting approaches for controlled knowledge
editing in LLMs, highlighting the need for methods that can operate within tight memory budgets.

\paragraph{Gap addressed by this work.}
No prior study directly compares rehearsal-based, regularization-based, and orthogonality-based
PEFT strategies for catastrophic forgetting mitigation under the strict compute budget of a single
commodity GPU (NVIDIA T4, 16\,GB) using 4-bit quantized small LLMs ($\leq$4B parameters) on a
standardized GLUE task sequence.
We fill this gap by benchmarking six methods---including our proposed Minimal Mixed Rehearsal
(MMR)---under identical hardware and hyperparameter constraints, providing actionable guidance for
practitioners operating in low-resource continual fine-tuning scenarios.
